% Options for packages loaded elsewhere
\PassOptionsToPackage{unicode}{hyperref}
\PassOptionsToPackage{hyphens}{url}
%
\documentclass[
  11pt,
]{article}
\usepackage{amsmath,amssymb}
\usepackage{iftex}
\ifPDFTeX
  \usepackage[T1]{fontenc}
  \usepackage[utf8]{inputenc}
  \usepackage{textcomp} % provide euro and other symbols
\else % if luatex or xetex
  \usepackage{unicode-math} % this also loads fontspec
  \defaultfontfeatures{Scale=MatchLowercase}
  \defaultfontfeatures[\rmfamily]{Ligatures=TeX,Scale=1}
\fi
\usepackage{lmodern}
\ifPDFTeX\else
  % xetex/luatex font selection
\fi
% Use upquote if available, for straight quotes in verbatim environments
\IfFileExists{upquote.sty}{\usepackage{upquote}}{}
\IfFileExists{microtype.sty}{% use microtype if available
  \usepackage[]{microtype}
  \UseMicrotypeSet[protrusion]{basicmath} % disable protrusion for tt fonts
}{}
\makeatletter
\@ifundefined{KOMAClassName}{% if non-KOMA class
  \IfFileExists{parskip.sty}{%
    \usepackage{parskip}
  }{% else
    \setlength{\parindent}{0pt}
    \setlength{\parskip}{6pt plus 2pt minus 1pt}}
}{% if KOMA class
  \KOMAoptions{parskip=half}}
\makeatother
\usepackage{xcolor}
\usepackage[margin=1in]{geometry}
\usepackage{color}
\usepackage{fancyvrb}
\newcommand{\VerbBar}{|}
\newcommand{\VERB}{\Verb[commandchars=\\\{\}]}
\DefineVerbatimEnvironment{Highlighting}{Verbatim}{commandchars=\\\{\}}
% Add ',fontsize=\small' for more characters per line
\usepackage{framed}
\definecolor{shadecolor}{RGB}{248,248,248}
\newenvironment{Shaded}{\begin{snugshade}}{\end{snugshade}}
\newcommand{\AlertTok}[1]{\textcolor[rgb]{0.94,0.16,0.16}{#1}}
\newcommand{\AnnotationTok}[1]{\textcolor[rgb]{0.56,0.35,0.01}{\textbf{\textit{#1}}}}
\newcommand{\AttributeTok}[1]{\textcolor[rgb]{0.13,0.29,0.53}{#1}}
\newcommand{\BaseNTok}[1]{\textcolor[rgb]{0.00,0.00,0.81}{#1}}
\newcommand{\BuiltInTok}[1]{#1}
\newcommand{\CharTok}[1]{\textcolor[rgb]{0.31,0.60,0.02}{#1}}
\newcommand{\CommentTok}[1]{\textcolor[rgb]{0.56,0.35,0.01}{\textit{#1}}}
\newcommand{\CommentVarTok}[1]{\textcolor[rgb]{0.56,0.35,0.01}{\textbf{\textit{#1}}}}
\newcommand{\ConstantTok}[1]{\textcolor[rgb]{0.56,0.35,0.01}{#1}}
\newcommand{\ControlFlowTok}[1]{\textcolor[rgb]{0.13,0.29,0.53}{\textbf{#1}}}
\newcommand{\DataTypeTok}[1]{\textcolor[rgb]{0.13,0.29,0.53}{#1}}
\newcommand{\DecValTok}[1]{\textcolor[rgb]{0.00,0.00,0.81}{#1}}
\newcommand{\DocumentationTok}[1]{\textcolor[rgb]{0.56,0.35,0.01}{\textbf{\textit{#1}}}}
\newcommand{\ErrorTok}[1]{\textcolor[rgb]{0.64,0.00,0.00}{\textbf{#1}}}
\newcommand{\ExtensionTok}[1]{#1}
\newcommand{\FloatTok}[1]{\textcolor[rgb]{0.00,0.00,0.81}{#1}}
\newcommand{\FunctionTok}[1]{\textcolor[rgb]{0.13,0.29,0.53}{\textbf{#1}}}
\newcommand{\ImportTok}[1]{#1}
\newcommand{\InformationTok}[1]{\textcolor[rgb]{0.56,0.35,0.01}{\textbf{\textit{#1}}}}
\newcommand{\KeywordTok}[1]{\textcolor[rgb]{0.13,0.29,0.53}{\textbf{#1}}}
\newcommand{\NormalTok}[1]{#1}
\newcommand{\OperatorTok}[1]{\textcolor[rgb]{0.81,0.36,0.00}{\textbf{#1}}}
\newcommand{\OtherTok}[1]{\textcolor[rgb]{0.56,0.35,0.01}{#1}}
\newcommand{\PreprocessorTok}[1]{\textcolor[rgb]{0.56,0.35,0.01}{\textit{#1}}}
\newcommand{\RegionMarkerTok}[1]{#1}
\newcommand{\SpecialCharTok}[1]{\textcolor[rgb]{0.81,0.36,0.00}{\textbf{#1}}}
\newcommand{\SpecialStringTok}[1]{\textcolor[rgb]{0.31,0.60,0.02}{#1}}
\newcommand{\StringTok}[1]{\textcolor[rgb]{0.31,0.60,0.02}{#1}}
\newcommand{\VariableTok}[1]{\textcolor[rgb]{0.00,0.00,0.00}{#1}}
\newcommand{\VerbatimStringTok}[1]{\textcolor[rgb]{0.31,0.60,0.02}{#1}}
\newcommand{\WarningTok}[1]{\textcolor[rgb]{0.56,0.35,0.01}{\textbf{\textit{#1}}}}
\setlength{\emergencystretch}{3em} % prevent overfull lines
\providecommand{\tightlist}{%
  \setlength{\itemsep}{0pt}\setlength{\parskip}{0pt}}
\setcounter{secnumdepth}{5}
\ifLuaTeX
  \usepackage{selnolig}  % disable illegal ligatures
\fi
\usepackage{bookmark}
\IfFileExists{xurl.sty}{\usepackage{xurl}}{} % add URL line breaks if available
\urlstyle{same}
\hypersetup{
  hidelinks,
  pdfcreator={LaTeX via pandoc}}

\author{}
\date{}

\begin{document}

{
\setcounter{tocdepth}{3}
\tableofcontents
}
\section{Project 2 - Unit Root
Testing}\label{project-2---unit-root-testing}

\textbf{Course:} Financial Econometrics\\
\textbf{Author:} Omid Karami Chamgordani - 400203402\\
\textbf{Date:} 2022-12-10

\subsection{Objective}\label{objective}

This report studies the time-series properties of an annual Brazilian
economic series retrieved from Quandl/BCB code \texttt{BCB/226}. The
main goal is to check whether the series has a unit root and whether
differencing is required before modeling.

\subsection{Data import and initial
plot}\label{data-import-and-initial-plot}

The data are imported at annual frequency and plotted over time. The
visual pattern suggests that the level series is non-stationary. For a
public GitHub repository, the API key is read from the environment
variable \texttt{QUANDL\_API\_KEY} rather than hard-coded in the script.

\begin{Shaded}
\begin{Highlighting}[]
\FunctionTok{library}\NormalTok{(Quandl)}
\FunctionTok{library}\NormalTok{(urca)}
\FunctionTok{library}\NormalTok{(tidyverse)}

\FunctionTok{Quandl.api\_key}\NormalTok{(}\FunctionTok{Sys.getenv}\NormalTok{(}\StringTok{"QUANDL\_API\_KEY"}\NormalTok{))}
\NormalTok{brazil\_c }\OtherTok{=} \FunctionTok{Quandl}\NormalTok{(}\StringTok{"BCB/226"}\NormalTok{, }\AttributeTok{collapse =} \StringTok{"annual"}\NormalTok{)}
\FunctionTok{plot}\NormalTok{(brazil\_c[,}\DecValTok{1}\NormalTok{], }\FunctionTok{ts}\NormalTok{(brazil\_c[,}\DecValTok{2}\NormalTok{]), }\AttributeTok{col=}\StringTok{"blue"}\NormalTok{, }\AttributeTok{lwd=}\DecValTok{2}\NormalTok{,}
     \AttributeTok{xlab=}\StringTok{"date"}\NormalTok{, }\AttributeTok{ylab=}\StringTok{"value"}\NormalTok{)}

\NormalTok{x }\OtherTok{=} \FunctionTok{ts}\NormalTok{(brazil\_c[,}\DecValTok{2}\NormalTok{])}
\FunctionTok{lm}\NormalTok{(x[}\DecValTok{2}\SpecialCharTok{:}\DecValTok{33}\NormalTok{] }\SpecialCharTok{\textasciitilde{}}\NormalTok{ x[}\DecValTok{1}\SpecialCharTok{:}\DecValTok{32}\NormalTok{])}
\end{Highlighting}
\end{Shaded}

The simple autoregression gives an intercept of approximately 1.2162 and
a lag coefficient of approximately 0.8191. Because the lag coefficient
is close to one and the level plot shows persistence, formal unit-root
tests are needed.

\subsection{Augmented Dickey-Fuller test with
trend}\label{augmented-dickey-fuller-test-with-trend}

The most general specification includes an intercept and deterministic
trend. The reported ADF statistic for the lagged level term is about
-0.8285, while the 5 percent critical value for \texttt{tau3} is about
-3.50. Therefore, the unit-root null is not rejected in the trend
specification.

\begin{Shaded}
\begin{Highlighting}[]
\CommentTok{\# step 1: trend specification}
\NormalTok{unitroot\_A }\OtherTok{\textless{}{-}} \FunctionTok{ur.df}\NormalTok{(x, }\AttributeTok{type=}\FunctionTok{c}\NormalTok{(}\StringTok{"trend"}\NormalTok{), }\AttributeTok{selectlags=}\FunctionTok{c}\NormalTok{(}\StringTok{"AIC"}\NormalTok{))}
\FunctionTok{summary}\NormalTok{(unitroot\_A)}
\FunctionTok{plot}\NormalTok{(unitroot\_A)}
\NormalTok{unitroot\_A}\SpecialCharTok{@}\NormalTok{teststat}
\end{Highlighting}
\end{Shaded}

\subsection{ADF test with drift}\label{adf-test-with-drift}

After removing the deterministic trend, the drift specification is
estimated. The ADF statistic is about -0.0529, which is again far above
the 5 percent critical value in absolute terms. The null of a unit root
is still not rejected.

\begin{Shaded}
\begin{Highlighting}[]
\CommentTok{\# step 5: drift specification}
\NormalTok{unitroot\_A }\OtherTok{\textless{}{-}} \FunctionTok{ur.df}\NormalTok{(x, }\AttributeTok{type=}\FunctionTok{c}\NormalTok{(}\StringTok{"drift"}\NormalTok{), }\AttributeTok{selectlags=}\FunctionTok{c}\NormalTok{(}\StringTok{"AIC"}\NormalTok{))}
\FunctionTok{summary}\NormalTok{(unitroot\_A)}
\FunctionTok{plot}\NormalTok{(unitroot\_A)}
\NormalTok{unitroot\_A}\SpecialCharTok{@}\NormalTok{teststat}
\end{Highlighting}
\end{Shaded}

\subsection{ADF test without deterministic
terms}\label{adf-test-without-deterministic-terms}

The final level specification removes both the trend and the intercept.
The test again supports the same conclusion: the level series behaves as
a non-stationary series.

\begin{Shaded}
\begin{Highlighting}[]
\CommentTok{\# step 8: no deterministic terms}
\NormalTok{unitroot\_A }\OtherTok{\textless{}{-}} \FunctionTok{ur.df}\NormalTok{(x, }\AttributeTok{type=}\FunctionTok{c}\NormalTok{(}\StringTok{"none"}\NormalTok{), }\AttributeTok{selectlags=}\FunctionTok{c}\NormalTok{(}\StringTok{"AIC"}\NormalTok{))}
\FunctionTok{summary}\NormalTok{(unitroot\_A)}
\FunctionTok{plot}\NormalTok{(unitroot\_A)}
\NormalTok{unitroot\_A}\SpecialCharTok{@}\NormalTok{teststat}
\end{Highlighting}
\end{Shaded}

\subsection{First difference}\label{first-difference}

The first difference of the series is then tested. After differencing,
the evidence supports stationarity more strongly than in the level
series. This suggests that the original series can be treated as
integrated of order one, I(1), for subsequent modeling.

\begin{Shaded}
\begin{Highlighting}[]
\CommentTok{\# First difference}
\NormalTok{dx }\OtherTok{\textless{}{-}} \FunctionTok{diff}\NormalTok{(x)}
\FunctionTok{plot}\NormalTok{(dx)}

\CommentTok{\# ADF on first difference}
\NormalTok{unitroot\_A }\OtherTok{\textless{}{-}} \FunctionTok{ur.df}\NormalTok{(dx, }\AttributeTok{type=}\FunctionTok{c}\NormalTok{(}\StringTok{"trend"}\NormalTok{), }\AttributeTok{selectlags=}\FunctionTok{c}\NormalTok{(}\StringTok{"AIC"}\NormalTok{))}
\FunctionTok{summary}\NormalTok{(unitroot\_A)}
\FunctionTok{plot}\NormalTok{(unitroot\_A)}
\NormalTok{unitroot\_A}\SpecialCharTok{@}\NormalTok{teststat}
\end{Highlighting}
\end{Shaded}

\subsection{Conclusion}\label{conclusion}

The level series is non-stationary under the common ADF specifications.
After first differencing, the series is more consistent with
stationarity. In future work, models should use the differenced series
or include an integration component.

\end{document}
